\documentclass[12pt]{article}
\usepackage[margin=0.6in,top=0.65in,bottom=0.5in]{geometry}
\usepackage{lmodern}
\usepackage{microtype}
\usepackage{fancyhdr}
\usepackage{titlesec}
\usepackage{enumitem}
\usepackage{lastpage}
\usepackage[hidelinks]{hyperref}

\setlength{\parindent}{0pt}
\setlength{\parskip}{0.22em}
\setlength{\headheight}{26pt}
\raggedbottom
\titleformat{\section}{\normalsize\bfseries\scshape}{}{0pt}{}
\titleformat{\subsection}{\normalsize\bfseries}{}{0pt}{}
\titlespacing*{\section}{0pt}{0.4em}{0.15em}
\titlespacing*{\subsection}{0pt}{0.3em}{0.1em}
\pagestyle{fancy}
\fancyhf{}
\fancyhead[C]{\footnotesize Lit Example Reader \quad Assignment ID: U1L7LE \quad Page \thepage\ of \pageref{LastPage}}
\renewcommand{\headrulewidth}{0.5pt}

\newenvironment{playpassage}{\begin{quote}\small\setlength{\parskip}{0.12em}\setlength{\parindent}{0pt}}{\end{quote}}
\newcommand{\speaker}[1]{\textbf{#1.}}

\begin{document}

\begin{center}
{\LARGE\bfseries Week 7 Lit Example Reader}\par\vspace{0.02em}
{\large\scshape Macbeth: Appearance of Proof and Pressure}\par\vspace{0.08em}
\rule{0.78\textwidth}{0.5pt}
\end{center}

\textbf{Purpose:} Macbeth passages for Week 7: diagnose reasoning that feels strong because of emotional or rhetorical pressure rather than earned proof.

\textbf{What to watch for:} What conclusion is being forced? What makes it feel proved? Where does the support fail?

\section*{Before the Passages: The Week 7 Logic Tool}

\begin{itemize}[leftmargin=1.3em,itemsep=0.08em,topsep=0.1em,parsep=0.04em]
\item \textbf{Fallacy:} reasoning that seems to prove a conclusion but fails to prove it.
\item \textbf{Appearance of proof:} confidence, shame, loyalty language, or vivid force that makes a claim feel finished.
\item \textbf{False claim vs bad inference:} untrue statement versus a leap that does not follow.
\item \textbf{Fallacy autopsy:} conclusion, apparent proof, defect, repair or narrowed claim.
\end{itemize}

\section*{Passage 1: Act 1, Scene 7 --- Lady Macbeth Presses Macbeth}

\textbf{Context:} Macbeth hesitates about killing Duncan. Lady Macbeth answers with a fierce attack on his resolve and manhood rather than a careful moral proof.

\begin{playpassage}
\speaker{Lady Macbeth} Was the hope drunk\\
Wherein you dress'd yourself? hath it slept since?\\
And wakes it now, to look so green and pale\\
At what it did so freely? From this time\\
Such I account thy love. Art thou afeard\\
To be the same in thine own act and valour\\
As thou art in desire? Wouldst thou have that\\
Which thou esteem'st the ornament of life,\\
And live a coward in thine own esteem,\\
Letting ``I dare not'' wait upon ``I would,''\\
Like the poor cat i' the adage?\\
\speaker{Macbeth} Prithee, peace:\\
I dare do all that may become a man;\\
Who dares do more is none.\\
\speaker{Lady Macbeth} What beast was't, then,\\
That made you break this enterprise to me?\\
When you durst do it, then you were a man;\\
And, to be more than what you were, you would\\
Be so much more the man.
\end{playpassage}

\subsection*{Reader Notes}

Here is your most assisted example. Lady Macbeth's pressure can be reconstructed as: a real man carries out the enterprise he named; you are hesitating; therefore you are less than a man / a coward in your own esteem. The speech feels like proof because it is intimate, shaming, and quick. But \textit{man} and \textit{enterprise} are loaded terms, and the leap from hesitation to proved unworthiness is not a demonstration that Duncan must die. Emotional force is doing work that logical force has not done.

\textbf{Reading task:} Underline the conclusion she presses. Circle shame or manhood-pressure words. In the margin, write one phrase that looks like proof and one place the leap may fail.

\newpage

\section*{Passage 2: Act 1, Scene 7 --- Macbeth Settles}

\textbf{Context:} After Lady Macbeth's pressure and practical plan, Macbeth's resistance collapses. Watch whether he has gained new proof or mainly accepted the pressure.

\begin{playpassage}
\speaker{Lady Macbeth} When Duncan is asleep---\\
his two chamberlains\\
Will I with wine and wassail so convince\\
That memory, the warder of the brain,\\
Shall be a fume\ldots\\
What cannot you and I perform upon\\
The unguarded Duncan? what not put upon\\
His spongy officers, who shall bear the guilt\\
Of our great quell?\\
\speaker{Macbeth} Will it not be received,\\
When we have mark'd with blood those sleepy two\\
Of his own chamber and used their very daggers,\\
That they have done't?\\
\speaker{Lady Macbeth} Who dares receive it other,\\
As we shall make our griefs and clamour roar\\
Upon his death?\\
\speaker{Macbeth} I am settled, and bend up\\
Each corporal agent to this terrible feat.\\
Away, and mock the time with fairest show:\\
False face must hide what the false heart doth know.
\end{playpassage}

\subsection*{Reader Notes}

This passage shifts more work to you. Lady Macbeth offers a plan for covering the murder, and Macbeth says he is ``settled.'' Mark what kind of support the plan actually supplies. Does a workable cover-up prove that the killing is justified, or does it only make the crime feel more possible?

\textbf{Reading task:} Number the practical steps in the plan. Beside ``I am settled,'' write whether the settlement rests on new proof that the murder is right, on confidence that they can hide it, or on both. Save the full autopsy for the worksheet.

\section*{How to Use These Passages on the Worksheet}

The worksheet asks for a fallacy autopsy, not a plot summary. Keep separate: what the speakers want, what their language makes feel necessary, and what has actually been proved.

\end{document}
